\documentclass[compress]{beamer}
\input{../../resources/preamble.sty}
\addbibresource{../../resources/literature.bib}
\graphicspath{{../../resources/img/}}


\begin{document}

\title[Big Data and Automated Content Analysis]{\textbf{Big Data and Automated Content Analysis (6EC)} 
\\Week 4: »Processing textual data«
\\Thursday}
\author[Anne Kroon]{Anne Kroon\\ \footnotesize{a.c.kroon@uva.nl}}
\date{April 24, 2025}
\institute[UvA CW]{UvA RM Communication Science}


\begin{frame}
	\titlepage
\end{frame}

\begin{frame}{Today}
	\tableofcontents
\end{frame}

\begin{frame}[standout]
Hopefully, everything is clear from last week, and that the take-home exam is going well.
\end{frame}


\begin{frame}[standout]
This week, we will get a general overview of working with textual data. This knowledge will help you to get started with cool automated content analyses techniques -- which we will start with next week.
\end{frame}


\input{../../modules/working-with-text/basic-string-operations.tex}

\input{../../modules/working-with-text/regular-expressions.tex}

\input{../../modules/working-with-text/bow.tex}



%\question{Any questions?}

\section{Next steps}

\begin{frame}[standout]
Make sure you understood all of today's concepts.

Re-read the chapters.

I prepared exercises to work on (alone or in teams):
\large{\url{https://github.com/uvacw/teaching-bdaca/blob/main/6ec-course/week04/exercises/}}
\\

\end{frame}


\begin{frame}[allowframebreaks,plain]
	\printbibliography
\end{frame}



\end{document}
